% Generated draft: requires mathematical and source review.
\documentclass[11pt]{article}
\usepackage[margin=25mm]{geometry}
\usepackage{fontspec}
\setmainfont{Latin Modern Roman}
\usepackage{amsmath,amssymb,mathtools,unicode-math}
\setmathfont{Latin Modern Math}
\usepackage{enumitem,needspace,xurl,hyperref}
\setlength{\parindent}{0pt}
\setlength{\parskip}{5pt}
\setlength{\emergencystretch}{3em}
\newcommand{\N}{\mathbb{N}}
\newcommand{\Z}{\mathbb{Z}}
\newcommand{\Q}{\mathbb{Q}}
\newcommand{\R}{\mathbb{R}}
\newcommand{\C}{\mathbb{C}}
\newcommand{\F}{\mathbb{F}}
\newcommand{\A}{\mathbb{A}}
\newcommand{\PP}{\mathbb{P}}
\newcommand{\E}{\mathbb{E}}
\newcommand{\eps}{\varepsilon}
\newcommand{\dd}{\,\mathrm{d}}
\begin{document}
\section*{UnsolvedMath ID: 30005327}
\textbf{Source problem code:} catalogue entry
\subsection{Definitions and mathematical statement}
For a positive integer \(m\), a permutation \(\pi\) of \(\{1,\ldots,m\}\) avoids the pattern \(321\) if there are no indices \(1\leq i<j<k\leq m\) with \(\pi(i)>\pi(j)>\pi(k)\).

The supplied record describes \(P_{321}(n)\) as counting parity-restricted \(321\)-avoiding permutations, but does not define the parity restriction, the underlying permutation sets, or their dependence on \(n\). Consequently, the sequence \(P_{321}(n)\) cannot be reconstructed from this record. Retaining its undefined notation, set formally
\[
F(q)=\sum_{n\geq 1}P_{321}(n)q^n.
\]
Here \(q\) is a formal indeterminate; the displayed series has no constant term.

To specify the properties being asked about, let \(K\) be a field of characteristic zero containing the intended coefficients and suppose \(F\in K[[q]]\). The source does not specify the intended coefficient field. The series \(F\) is algebraic over \(K(q)\) if there exists a nonzero polynomial \(A(q,Y)\in K[q,Y]\) such that \(A(q,F(q))=0\). It is D-finite over \(K(q)\) if there exist an integer \(r\geq 0\) and polynomials \(p_0(q),\ldots,p_r(q)\in K[q]\), not all zero, such that
\[
\sum_{j=0}^{r}p_j(q)F^{(j)}(q)=0.
\]
It is D-algebraic over \(K(q)\) if there exist an integer \(r\geq 0\) and a nonzero polynomial \(B(q,Y_0,\ldots,Y_r)\in K[q,Y_0,\ldots,Y_r]\) such that
\[
B\bigl(q,F(q),F'(q),\ldots,F^{(r)}(q)\bigr)=0.
\]
All derivatives here are formal derivatives, with \(F^{(0)}=F\).

The source asks whether \(F\) has each of these three properties and asks for the asymptotic behavior of \(P_{321}(n)\) as \(n\to\infty\). These questions remain conditional on recovering the definition of \(P_{321}(n)\) and specifying the intended coefficient field. The supplied record labels the problem open; this draft does not independently establish that status.
\subsection{Short English statement}
Determine whether the generating function for the intended parity-restricted \(321\)-avoiding permutations is algebraic, D-finite, or D-algebraic, and determine the asymptotic behavior of its coefficients. The supplied record omits the parity restriction and the definition of \(P_{321}(n)\), so the counting problem requires clarification before a complete mathematical statement is possible.
\subsection{Sources}
\begin{itemize}
\item Supplied local registry record, permanent UnsolvedMath ID \(30005327\), source problem code \texttt{catalogue entry}, titled ``Generating Functions for Parity Restricted 321 Avoiding Permutations.'' The record gives proposed year 2022 and supplies no named proposer.
\item Supplied report citation: \emph{Enumerative Combinatorics (2023)}, Oberwolfach Reports, DOI 10.4171/owr/2022/57, pp.\ 3215--3305. Original OWR report: \url{https://doi.org/10.4171/owr/2022/57}.
\item Supplied literature source labeled ``problem-session paper'': \url{https://personal.utdallas.edu/~nxw170830/docs/Papers/Misc/Oberwolfach_EC_2022_problem_session.pdf}.
\item Literature assessment dated exactly \texttt{2026-08-22}: the supplied curation reports that the algebraicity, D-finiteness, D-algebraicity, and coefficient asymptotics remain unresolved in the identified literature. This is supplied dated curation, not an independently verified assessment; the record explicitly calls for independent verification of time-sensitive status and cited claims.
\item The registry labels its statement \texttt{corrected\_verified}, explaining that only formula delimiters, notation, grammar, and typography were repaired. That supplied label does not resolve the missing definition or constitute independent mathematical verification in this draft.
\end{itemize}
\end{document}
